A separate sensitivity test perturbs the supplied poses while keeping evaluation geometry and world anchors fixed. For each frame, independent Gaussian translation perturbations have per-axis standard deviation 0, 1, or 3\,cm; rotation-vector components have corresponding standard deviations 0, 1, or 3 degrees. Rotation is applied about the camera center. The same standardized perturbations are reused across severity levels for each seed. This tests artificial pose jitter, not the trajectory drift of a real SLAM system.

Across three seeds at $B=256$, mean accumulated-surface error on \texttt{desk} rises from \DeskPoseZeroSeen{}\,cm with oracle poses to \DeskPoseOneSeen{} and \DeskPoseThreeSeen{}\,cm at the two nonzero levels. On \texttt{xyz}, it rises from \XyzPoseZeroSeen{} to \XyzPoseOneSeen{} and \XyzPoseThreeSeen{}\,cm. Focal error rises from \DeskPoseZeroFocus{} to \DeskPoseOneFocus{} and \DeskPoseThreeFocus{}\,cm on \texttt{desk}, and from \XyzPoseZeroFocus{} to \XyzPoseOneFocus{} and \XyzPoseThreeFocus{}\,cm on \texttt{xyz}. The zero-noise runs exactly reproduce the corresponding primary accuracy summaries. Memory preserves supplied geometry; it cannot correct erroneous registration.

\begin{table}[t]
\centering\small
\caption{Additional rotation sequence, fr1/rpy, at $B=256$. Mean errors in cm across three seeds; same fixed protocol. This sequence was added after the two-sequence primary protocol was fixed.}
\label{tab:rotation}
\begin{tabular}{lrr}
\toprule Method & All observed & Focal \\
\midrule
Unregistered F-GNG & 103.93 & 52.60 \\
World-only F-GNG & 106.45 & 57.86 \\
\textbf{WA-FGNG} & 14.68 & 6.89 \\
Replay, uniform field & 14.39 & 9.11 \\
Replay + DBL-GNG & 13.44 & 9.03 \\
Replay + FPS & 15.10 & 14.36 \\
\bottomrule\end{tabular}
\end{table}

An additional fr1/rpy experiment uses the same parameters, node limit 256, and three seeds, yielding 18 further runs. Its 120 selected frames span 24.1\,s. As Table~\ref{tab:rotation} shows, foveation reduces focal error from \RpyUniformFocus{} to \RpyMemoryFocus{}\,cm relative to the matched uniform-field ablation, a \RpyFocusGain\% improvement, with a \RpyGlobalCost\% increase in global error. Unsupported-edge fraction remains \RpyMemoryEdge\% for \method{}. This adds evidence under rotational camera motion, while remaining within the benchmark's narrow office/sensor regime.
